% Separately instantiated corrected French for printed p.120.
% Decisions: TOH-P120-R1-SRC-DISPOSITION-001; TOH-P120-R2-REV-001.
% No authorial error is established; the added target is zero-width metadata.
\SourcePage{120}

\LayerTarget{tohoku:introduction:I:p4:cont:p120}
que la théorie s'applique aussi aux espaces non séparés qui interviennent en
Géométrie Algébrique abstraite ou en ``Géométrie Arithmétique'' [15] [8]. Des
conversations avec R. Godement et H. Cartan ont été très précieuses pour la
mise au point de la théorie, et en particulier l'introduction par Godement des
\emph{faisceaux flasques} et des \emph{faisceaux mous}, qui se substituent
avantageusement aux faisceaux fins dans bien des questions, s'est révélée
extrêmement commode. Un exposé plus complet, auquel nous renverrons pour divers
points de détail, sera donné dans un livre en préparation par R. Godement [9].

\LayerTarget{tohoku:introduction:I:p5}
Le Chapitre IV traite la question non classique des Ext de faisceaux de
modules, on y trouvera en particulier une suite spectrale utile qui relie les
Ext ``globaux'' et les Ext ``locaux''.
\LayerTarget{tohoku:introduction:I:p6}
La situation se corse au Chapitre V, où de plus un groupe $G$ opère sur
l'espace $X$, le faisceau d'anneaux $\mathbf{O}$ donné sur $X$, et les faisceaux de
modules sur $\mathbf{O}$ qu'on considère. On obtient en particulier dans 5.2 un énoncé
qui me semble être la forme définitive de la théorie cohomologique
``Čechiste'' des espaces à groupe (non topologique) d'opérateurs, pouvant avoir
des points fixes. Il s'exprime en introduisant de nouveaux foncteurs
$H^n(X;G,A)$ (implicites déjà dans bien des cas particuliers antérieurs): on
trouve alors deux foncteurs spectraux, à termes initiaux remarquables, qui y
aboutissent.

\LayerTarget{tohoku:introduction:II}
\LayerContents{II. Applications}
\paragraph{II. Applications.}
\LayerTarget{tohoku:introduction:II:p1}
Faute de place, je n'ai pu donner dans cet exposé que très peu d'applications
des techniques employées (notamment dans 3.4 et 3.6), me contentant d'en
signaler quelques unes au passage. Signalons encore les applications
suivantes:

\LayerTarget{tohoku:introduction:II:a}
\par a)
La notion de Ext de faisceaux de modules permet la formulation la plus
générale connue du ``théorème de dualité algébrique'' de Serre: Si $A$ est un
faisceau algébrique cohérent [15] sur une variété algébrique projective de
dimension $n$ sans singularités, alors le dual de $H^p(X,A)$ s'identifie
canoniquement à $\mathrm{Ext}_{\mathrm{O}}^{n-p}(X;A,\Omega^n)$, où $\mathbf{O}$ (resp.
$\Omega^n$) est le faisceau des germes de fonctions régulières (resp. de
$n$-formes régulières) sur $X$.

\LayerTarget{tohoku:introduction:II:b}
\par b)
Tout le formalisme développé dans les Chapitres III, IV, V peut s'appliquer en
Géométrie Algébrique Abstraite. Je montrerai ailleurs comment il permet
d'étendre aux variétés algébriques complètes divers résultats prouvés par
Serre [15] [16] [17] pour les variétés projectives.

\LayerTarget{tohoku:introduction:II:c}
\par c)
Il semble que les $H^n(X;G,A)$ soient l'intermédiaire naturel pour une théorie
générale des puissances réduites de Steenrod dans les faisceaux, et de la
cohomologie des puissances symétriques d'espaces quelconques, théorie qui
s'applique aussi en Géométrie Algébrique en car. $p$.

\LayerTarget{tohoku:introduction:III}
\LayerContents{III. Lacunes}
\paragraph{III. Lacunes.}
\LayerTarget{tohoku:introduction:III:p1}
Pour ne pas allonger cet exposé, j'ai passé sous silence les questions de
structures multiplicatives, quoiqu'elles soient tout à fait essentielles dans
l'application des notions des Chapitres III, IV, V. Signalons d'ailleurs qu'il
ne semble y avoir encore de théorie satisfaisante des structures
multiplicatives en Algèbre Homologique, ayant le degré de généralité et de
simplicité nécessaire (le Chapitre II de [6] étant d'ailleurs une illustration
frappante de cet état de choses).\LayerTarget{tohoku:introduction:footnote:1bis}\textsuperscript{1 bis)} Pour la multiplication
en cohomologie

\TohokuContinuation{121}
